\documentclass[preprint, 3p,
authoryear]{elsarticle} %review=doublespace preprint=single 5p=2 column
%%% Begin My package additions %%%%%%%%%%%%%%%%%%%

\usepackage[hyphens]{url}

  \journal{CERC 2026} % Sets Journal name

\usepackage{graphicx}
%%%%%%%%%%%%%%%% end my additions to header

\usepackage[T1]{fontenc}
\usepackage{lmodern}
\usepackage{amssymb,amsmath}
% TODO: Currently lineno needs to be loaded after amsmath because of conflict
% https://github.com/latex-lineno/lineno/issues/5
\usepackage{lineno} % add
\usepackage{ifxetex,ifluatex}
\usepackage{fixltx2e} % provides \textsubscript
% use upquote if available, for straight quotes in verbatim environments
\IfFileExists{upquote.sty}{\usepackage{upquote}}{}
\ifnum 0\ifxetex 1\fi\ifluatex 1\fi=0 % if pdftex
  \usepackage[utf8]{inputenc}
\else % if luatex or xelatex
  \usepackage{fontspec}
  \ifxetex
    \usepackage{xltxtra,xunicode}
  \fi
  \defaultfontfeatures{Mapping=tex-text,Scale=MatchLowercase}
  \newcommand{\euro}{€}
\fi
% use microtype if available
\IfFileExists{microtype.sty}{\usepackage{microtype}}{}
\usepackage[]{natbib}
\bibliographystyle{plainnat}

\ifxetex
  \usepackage[setpagesize=false, % page size defined by xetex
              unicode=false, % unicode breaks when used with xetex
              xetex]{hyperref}
\else
  \usepackage[unicode=true]{hyperref}
\fi
\hypersetup{breaklinks=true,
            bookmarks=true,
            pdfauthor={},
            pdftitle={Do-It-Yourself information provision with mobilityHubTools},
            colorlinks=false,
            urlcolor=blue,
            linkcolor=magenta,
            pdfborder={0 0 0}}

\setcounter{secnumdepth}{5}
% Pandoc toggle for numbering sections (defaults to be off)


% tightlist command for lists without linebreak
\providecommand{\tightlist}{%
  \setlength{\itemsep}{0pt}\setlength{\parskip}{0pt}}




\usepackage{setspace}
\doublespacing
\usepackage{subfig}
\usepackage{booktabs}
\usepackage{longtable}
\usepackage{array}
\usepackage{multirow}
\usepackage{wrapfig}
\usepackage{float}
\usepackage{colortbl}
\usepackage{pdflscape}
\usepackage{tabu}
\usepackage{threeparttable}
\usepackage{threeparttablex}
\usepackage[normalem]{ulem}
\usepackage{makecell}
\usepackage{xcolor}



\begin{document}


\begin{frontmatter}

  \title{Do-It-Yourself information provision with mobilityHubTools}
    \author[]{James Reynolds%
  \corref{cor1}%
  \fnref{1}}
   \ead{julianj.reynolds@atu.ie} 
    \author[]{Brian McCann%
  %
  \fnref{1}}
   \ead{brian.mccann@atu.ie} 
    \author[]{Noreen Armstrong%
  %
  \fnref{2}}
   \ead{noreen.armstrong@atu.ie} 
    \author[]{Sarbast Molsem%
  %
  \fnref{3}}
   \ead{molsem@tcd.ie} 
    \author[]{Agnieszka Stefaniec%
  %
  \fnref{4}}
   \ead{a.stefaniec@soton.ac.uk} 
    \author[]{Brian Caulfield%
  %
  \fnref{3}}
   \ead{brian.caulfield@tcd.ie} 
    \author[]{Amaya Vega%
  %
  \fnref{2}}
   \ead{amaya.Vega@atu.ie} 
      \cortext[cor1]{Corresponding author}
    \fntext[1]{Atlantic Technological University, Sligo campus, Ireland}
    \fntext[2]{Atlantic Technological University, Galway City campus,
Ireland}
    \fntext[3]{Centre for Transport Research, Dept of Civil, Structural
and Environmental Engineering, Trinity College Dublin, Ireland}
    \fntext[4]{Department of Decision Analytics and Risk at Southampton
Business School, University of Southampton, UK}
  
  \begin{abstract}
  Mobility hubs are places where shared bikes/e-bikes, cars, e-scooters
  or similar modes are co-location, often close to transit and
  facilities such as charging stations or parcel lockers. The aim is to
  provide a combined alternative to (petrol-driven) car ownership and to
  support more sustainable travel behaviours. Signage, innformation
  displays an user guidance will likely be important at such hubs, but
  not much is known about industry best practices. There do not appear
  to be tools available to help mobility hub implementers, including for
  creating hubs where existing services are already co-located. This
  paper reports case studies of existing practices, and the development
  of software tools for creating maps and informational displays for
  mobility hubs in any location. Are there multiple mobility offerings
  in close proximity in your community? These tools, a printer and
  somewhere hang a poster might be all that is needed to spread the
  word.
  \end{abstract}
    \begin{keyword}
    Mobility hub, informational displays, shared modes, electric
mobility \sep 
    Mobility hub, informational displays, shared modes, electric
mobility
  \end{keyword}
  
 \end{frontmatter}

\section{Introduction}\label{introduction}

Shared mobility hubs are places where a range of modes and other
facilities may be provided to support more sustainable travel. They
often include shared, bikes, cargobikes, e-bikes, e-scooters or other
micro-mobility vehicles, as well as car share, public car charging and
other services. Together, these might help support shifts towards more
walking, cycling and transit use, and otherwise improve the
sustainability of transport systems through providing options for those
(more occasional) trips that are difficult without a private car.

Many mobility hubs have already been implemented, but with signage, maps
and other informational material created by individual agencies on a
site-by-site basis. Current and best practices for information provision
and user guidance at mobility hubs are unclear. As well, the graphic
design, mapping and similar effort needed to develop such material for
each site may limit the ability for practitioners to keep information up
to date, to trial new hubs, or to use bottom-up implementation
approaches to create hubs where shared or other modes are already in
close proximity.

This paper aims to explore and better understand mobility hub
information provision, and to develop tools to support practitioners. It
reports findings from case research examining existing practices at a
selection of mobility hubs. Also reported are findings from design-based
research that develops software for producing maps and information
displays at other mobility hubs, based on both replication and synthesis
of industry practices. The rest of this paper is structured as follows:
the next section outlines the context for this study. The research
methodology is then described, followed by presentation of the results.
These are then briefly discussed in a concluding section that also
identifies directions for future research.

\section{Research context}\label{research-context}

An early version of the mobility hub concept was developed in Toronto's
``Big Move'' transport plan \citep{Metrolinx:2008aa} in the late 2000s.
There the term was used to designate selected major regional transit
stitions around which landuses would be intensified and where intermodal
transfers would be facilitated
\citep[\citet{Engel-Yan:2012aa}]{Salsberg:2010aa}. Since then, the
mobility hub concept has expanded towards co-locating shared and/or
micro-mobility services often, but not always, close to transit. They
might also incorporate facilities and ``should be considered not just a
transfer node but as a place where a variety of services come together,
such as logistics services or restaurants'' \citep{Roukouni:2023aa}.
This definition from \citet{Roukouni:2023aa} is based on a review of the
literature, and also highlights the need for ``integration'',
``seamless, flexible connection'' and ``\ldots an enjoyable experience
for travelers'' . How this might be achieved, implemented and maintained
by practitioners remains unclear, but signs, maps and information
displays may be an important tool.

More broadly, `Movement and Place' and `Complete Streets'
\citep{Delbosc:2018aa}, and tactical urbanism approaches
\citep{Lydon:2015aa} suggest way to challenge prevailing naratives about
road space allocation, give communities greater power over (their own)
urban environments, and demonstrate the potential of change. Bottom-up
and incremental approaches, `pop-ups' and more formal trials have also
been found to be pragmatic strategies for building legitimacy
\emph{through} implementation \citep{Reynolds:2021ab}. In the context of
mobility hubs these might rely on rapid prototyping, iterative co-design
or similar approaches and a need for practitioners, tactical urbanists
and others to need adopt a\\
`do it yourself' attitude to developing information displays and
supporting maps.

As of yet, there does not appear to be standarised approaches for
providing such information, or a clear understanding of best practices
for user guidance at mobility hubs. Nor are there tools that might
allow\\
such best practices to be quickly and easily adapted to local contexts,
and for updates to be made as a mobility hub or its surrounding urban
environment changes over time. The next section, therefore describes the
case- and design-based research methodology adopted in this study to (at
least start) filling this gap in knowledge.

\section{Methodology}\label{methodology}

Research described in this paper aims to improve understanding of `best
practices' in in mobility hub signage and information provision, and to
develop tools to support mobility hub implementation. The research
questions relate to: \emph{How} information is provided at and about
mobility hubs; \emph{How} information at and about mobility hubs might
ideally be provided; and \emph{how} such best practices can be made
generally applicable?\footnote{More formally, the (null) hypothesis
  tested in this study is that it is \emph{not} possible to develop
  software tools for generating mobility hub information displays that
  reflect best practices and are generally applicable to other sites.}

\subsection{Study design}\label{study-design}

Study design is based around the Double Diamond design process
\citep{UK-Design-Council:0aa}, which \citet{Coxon:2018aa} and
\citet{Coxon:2021ab} identify as a suitable tool for design-based
research seeking to develop solutions to mobility problems. The Double
Diamond process involves four phases: 1. \emph{Discovery}, 2.
\emph{Definition}, 3. \emph{Development} and 4. \emph{Delivery}. These
phases alternate between divergent (1 \& 3) and convergent (2 \& 4)
approaches, although the overall process is iterative and may involve
multiple passes through each phase.

\subsection{Participants, material and
instruments}\label{participants-material-and-instruments}

The authors were the only participants in the research reported in this
paper. The research described in this paper is part of the larger
Electric shaRed mOBility hUbS Trial (ROBUST) project
\citep{ROBUST:2025aa}. Further research as part of the ROBUST project is
expected to involve mobility hub users, stakeholders and others as study
participants for testing, co-design and continued iteration of
information displays and other hub elements. These, however, are beyond
the scope of this paper.

Material for the study was drawn from: photos taken during sites visits
made by the lead author in the \emph{Discovery} phase; and mapping,
imagery and other online resources. Open Street Map (OSM) datasets
provided an input for the \emph{Development} phase.

Instruments included synthesis tables, developed as per
\citet{Miles:1994aa}. The \emph{Development} phase used the R
programming language. Package development during the \emph{Delivery}
phase was guided by \citet{wickham2023r}.

\subsection{Procedure and analysis}\label{procedure-and-analysis}

The \emph{Discovery} phase involved case research, which
\citet{Denscombe:2007aa}, \citet{Yin:2018ab} and others identify as
being especially effective for when: theory development is at an early
stage; research questions relate to `how' or `why'; and when exploring
phenomena within real-world contexts (as is the situation here). Cases
were selected from
https://www.como.org.uk/mobility-hubs/built-andplanned-hubs. Inclusion
was somewhat constrained by logistics, but guided by the case research
methodological literature (see \citet{Yin:2018ab} and others) to include
the leading examples of the mobility hubs developed by the West of
England Mayoral Combined Authority (WEMCA, Bristol) and Transport for
West Midlands (Birmingham). More representative examples from Camden
Council (London) and the London Borough of Redbridge (Woodford and
Wanstead) were also included to provide polar opposites for cross-case
comparison\footnote{The unit of analysis was the implementing agency /
  location. Within each case a small number of mobility hubs were
  visited to allow for within-case comparison. These were:
  (\emph{Redbridge}) South Woodford and High Street, Wanstead;
  (\emph{Camden}) Goldington Crescent, and Doric Way; (\emph{Bristol})
  Abbey Wood, UWE Frenchay, and Gainsborough Square; and
  (\emph{Birmingham}) Huntingtree Park and Andrew Road}.

The \emph{Definition} phase involved synthesis of findings from the site
visits and about each included case. This provided the basis for
defining the problem to be addressed by the research outputs of the
\emph{Development} and \emph{Delivery} phases.

The \emph{Development} phase involved iterative software coding to
replicate the Bristol and Birmingham style of information displays,
first undertaken for cases in the \emph{Discovery} phase, and then
repeated to test the software on the context of the ROBUST hub in
Galway. Finally, a combined display style was developed based on the
synthesis of industry practices and to respond to the problem developed
in the \emph{Definition} phase. The research also included iteration
through all four phases\footnote{ These iterations were tracked using
  the Git version control system, and code used in the production of
  this paper is available at
  \url{https://github.com/James-Reynolds/WP5_mobilityHubTools_papers/blob/main/ReynoldsEtAl2026_CERC.Rmd}.
  The software package itself is available at
  \url{https://github.com/James-Reynolds/mobilityHubTools}, and is
  expected to be developed further during later parts of the ROBUST
  project}, but in the interests of clarity the next section presents
the results in a linear format, stepping through each phase in turn.

\section{Results}\label{results}

\subsection{Discovery}\label{discovery}

Table \ref{tab:synthesis_table} summarises findings from the case
research\footnote{Table notes: 1. Case study questions (column 1) are of
  the form `How are details related to \textless COLUMN 1\textgreater{}
  at the hub communicated?' 2. An `X' denotes an explicit mention of the
  feature in hub information. `N' denotes no explicit mention, and a
  blank cell indicates the feature is not present at the hub. `S'
  denotes secure bicycle parking, such as a locker. 3. Totem denotes a
  free standing information pedestal. 4. The totem information display
  at the Doric Way mobility hub is specific to the Transport For London
  `Cycles' docked bike share facility. There is no information provided
  specifically about the mobility hub at Doric Way.}. Figures
\ref{fig:selected_images} show selected images from various hubs.

\begingroup\fontsize{7}{9}\selectfont

\begin{longtable}[t]{>{\raggedright\arraybackslash}p{2cm}|>{\centering\arraybackslash}p{1.1cm}>{\centering\arraybackslash}p{1.1cm}|>{\centering\arraybackslash}p{1.1cm}>{\centering\arraybackslash}p{1.1cm}|>{\centering\arraybackslash}p{1.1cm}>{\centering\arraybackslash}p{1.1cm}>{\centering\arraybackslash}p{1.5cm}|>{\centering\arraybackslash}p{1.5cm}>{\centering\arraybackslash}p{1.5cm}|}
\caption{\label{tab:synthesis_table}Hub characteristics. Source: Authors' assessment.}\\
\toprule
\multicolumn{1}{c}{ } & \multicolumn{2}{c}{Borough of Redbridge} & \multicolumn{2}{c}{Camden Council} & \multicolumn{3}{c}{WEMCA} & \multicolumn{2}{c}{Transport for West Midlands} \\
\cmidrule(l{3pt}r{3pt}){2-3} \cmidrule(l{3pt}r{3pt}){4-5} \cmidrule(l{3pt}r{3pt}){6-8} \cmidrule(l{3pt}r{3pt}){9-10}
Question? & Sth Wfd & High St, Wanstead & Goldington Crescent & Doric Way & Abbey Wood & UWE & Gainsborough Square & Huntingtree Park & Andrew Road\\
\midrule
Location & London & London & London & London & Bristol & Bristol & Bristol & Birmingham & Birmingham\\
\midrule\\
Information & Placard & Placard & None & Totem & Totem & Totem & Totem & Totem, poster & Totem, poster\\
Name & Mobility hub & Mobility hub & None & None & None & None & None & Local Travel Point & Local Travel Point\\
\midrule\\
Local map &  &  &  & X & X & X & X & X & X\\
Surrounds map &  &  &  & X & X & X & X & X & X\\
\addlinespace
Linemarking & X & X & X & X & X & X & X &  & \\
Indicators &  &  &  & X & X & X & X &  & \\
\midrule\\
Train & X &  &  & X & N & N & N &  & \\
Bus & X & N &  & X & X & X & X & X & X\\
Taxi & X &  &  & X &  &  &  &  & \\
\midrule\\
\addlinespace
Car club &  &  &  & N & X &  &  &  & N\\
E-car club & X &  & N &  &  &  &  &  & N\\
Share bike &  &  & N & X &  &  &  & X & X\\
E-bike &  &  & N & N & X & X & X &  & \\
E-cargobike &  &  & N &  & X & X &  &  & \\
\addlinespace
E-scooter &  &  & N &  &  & X & X &  & \\
\midrule\\
Car charger & N & N & N &  &  &  &  &  & \\
Bike parking & N & N & N, S & N, S & X, S & X, S & X & N, X & X, S\\
Bike repair &  &  & N &  & X & X & X &  & \\
Seating & X & N &  &  & X &  &  &  & \\
\addlinespace
Phone charge &  &  &  &  & X & X & X &  & \\
Wifi &  &  &  &  &  & X & X &  & \\
Toilets &  &  &  &  &  & X & X &  & \\
Cafe & N &  &  &  & X & X &  &  & \\
Food / shops &  &  &  &  & X & X &  & X & X\\
\bottomrule
\end{longtable}
\endgroup{}

\begin{figure}

{\centering \includegraphics[width=1\linewidth]{graphics/Figure1} 

}

\caption{Selected hub images. Source: authors.}\label{fig:selected_images}
\end{figure}

Information displays ranged from small placard signs (South Woodford and
Wanstead) to more comprehensive totem displays and posters (Bristol and
Birmingham). In Camden there were no information signs specific to the
hubs. Sites were labelled as ``Mobility Hub'' or ``Village Hub'' (South
Woodford and Wanstead) and ``Local Travel Point'' in Birmingham. There
was no labelling on-site in Camden or Bristol. Where present, totems all
included two maps: one at a local scale (within 5 minutes walk) and a
second showing a larger area (within 15 minutes). These generally showed
nearby transit, cafes, shops and other points of interest. Text
descriptions of local walking and cycling routes, hub facilities and
transit services were also included. Many of the sites had linemarking
that highlighted some hub components. The Bristol sites also included
cube shaped indicators mounted on poles to highlight individual hub
elements (Figure 1, bottom left).

Bicycle parking hoops were the most commonly provided feature. Some hubs
in Camden, Bristol and Birmingham also included bicycle lockers or
similar higher security parking. Seating and nearby toilets were
available as part of some hubs, but none included comprehensive shelter
or access to water. The Bristol hubs included USB connections for phone
recharging, and wifi internet connectivity at some locations.

Overall, findings during the \emph{Discovery} phase suggest that current
practices relating to information displays at mobility hubs are
inconsistent. Best practices appear to be represented by the totem
displays in Bristol and Birmingham.

\subsection{Definition}\label{definition}

Findings from the case research, iteration through the Double Diamond
phases and synthesis resulted led to the expression of the design
problem to be addressed in this research in the form of: \emph{five
questions a mobility hub information display needs answer}. These are
shown in Table \ref{tab:definition_details}. Also shown is an assessment
of how each of the cases does (or does not) respond to each.

\begingroup\fontsize{7}{9}\selectfont

\begin{longtable}[t]{>{\raggedright\arraybackslash}p{2.25cm}|>{\raggedright\arraybackslash}p{2.8cm}|>{\raggedright\arraybackslash}p{2.8cm}|>{\raggedright\arraybackslash}p{1.75cm}|>{\raggedright\arraybackslash}p{2.75cm}|>{\raggedright\arraybackslash}p{2.75cm}|}
\caption{\label{tab:definition_details}Defining the design problem through five questions. Authors' concept.}\\
\toprule
\multicolumn{2}{c}{Problem} & \multicolumn{4}{c}{Answers from cases} \\
\cmidrule(l{3pt}r{3pt}){1-2} \cmidrule(l{3pt}r{3pt}){3-6}
Question & Purpose & Redbridge London & Camden London & WEMCA Bristol & TWM Birmingham\\
\midrule
1. What is this? & Define current location as a mobility hub & Mobility or village hub, or 'Public Space' & Not answered & Not explictly answered & Local Travel Point\\
2. Where am I? & Identify current location & Place name & Not answered & Place name & Place name\\
3. What is here? & Identify hub elements & Limited & Not answered & Icons, cube markers, description & Icons, description\\
4. What is nearby? & Identify nearby services/facilities & Limited & Not answered & Local and regional maps, description, direction arrows & Local and regional maps, description, direction arrows\\
5. How can I travel from here? & Identify spatial network context & Not answered & Not answered & Description of bus/bike, bike routes on area map & Description of bus routes only\\
\bottomrule
\end{longtable}
\endgroup{}

The first question shown in Table \ref{tab:definition_details} relates
to identifying a mobility hub as something in and of itself. Various
terms were used in practice (Village Hub, Public Space, Local Travel
Point), except in Camden (no information display) and Bristol (place
name only).

The second question (\emph{Where am I?}) was answered in the cases
either by place names (Redbridge, Bristol and Birmingham) or not at all
(Camden). For the third question (\emph{What is here?}) the best
industry practices appear to be identify hub elements through text
descriptions, icons, linemarking and cube markers. Notably, hub elements
were not always shown on the maps answering the fourth question
(\emph{What is nearby?}). Answers to the final question shown in Table
\ref{tab:definition_details} (\emph{How can I travel from here?}) were
similarly limited, although some information displays included text
descriptions and limited mapping.

\subsection{Development}\label{development}

\begin{figure}

{\centering \includegraphics[width=1\linewidth]{graphics/Figure2} 

}

\caption{Software output: Birmingham (left), Bristol (centre) and synthesis (right) display types. Source: authors.}\label{fig:display_outputs_figure}
\end{figure}

During the \emph{Development} phase the R language was used to to:
download Open Street Maps (OSM) data; produce local and regional maps;
replicate the Bristol and Birmingham styles of totems; and create a
\emph{synthesis} display format, which responds to the problem
\emph{Definition}. OSM and R were chosen to facilitate reuse of the
software, as these are both open source and freely available. Figure
\ref{fig:display_outputs_figure} shows outputs of the software generated
for the Gainsborough Square case location.\\
These are intended to be printed at A0-scale, but the inset in Figure
\ref{fig:display_outputs_figure} (bottom left) demonstrates how the
mapping labels points of interests The maps include iconographic
representation of nearby facilities such as bus stops and cafes, which
can be generated automatically from OSM data or from user inputs
(i.e.~lat, lon). Maps included in the synthesis (Figure
\ref{fig:display_outputs_figure}, right) leverage map tiles associated
with OSM, including the the ``Transport'' and ``Cycle'' tile sets from
\citet{Allan:0aa}. Higher resolution images of the software outputs and
totems for the Galway ROBUST hub are available on github, along with the
software used to prepare this paper\footnote{See
  https://github.com/James-Reynolds/WP5\_mobilityHubTools\_papers/tree/main/results}.

\subsection{Delivery}\label{delivery}

Software developed in this research was assembled into a package of
functions, as per the guidance provided by \citet{wickham2023r}. This
package is available at
\url{https://github.com/James-Reynolds/mobilityHubTools}, and includes
documentation and test data.

\section{Discussion and conclusions}\label{discussion-and-conclusions}

The cases studied during the first phase of this research suggest much
variability in information displays and user guidance at mobility hubs.
Practices include comprehensive displays in (Bristol and Birmingham) and
cube shaped markers identifying hub elements (Bristol).\\
At the other end of the spectrum, sites in Camden had only linemarking
and did not appear to be identifiable as linked facilities or integrated
hubs.

The design problem identified in this study relates specifically to
totems and mapping, but clearly there are opportunities for further
research into mobility hub wayfinding and information provision. As
well, even those more extensive information displays found at some sites
do not appear to have fully answered the five questions defined here as
the design problem to be addressed. While the synthesis totem display
output from the \emph{Development} phase of the research described here
does answer all five questions (multiple times), further design
improvements can clearly be made. There may be a need to better
understand what facilities and services might best be included as `part'
of a hub, as opposed to those that are mapped as being `nearby'. More
immediately, further development of the software and user testing of the
software outputs would seem a direction for future research. As briefly
noted earlier, the research described in this paper is part of the
ROBUST project, which involves trials of shared electric mobility hubs
in various locations across Ireland. Part of this larger program of
research will include looking at user behaviour, usability and related
issues under real-world conditions. Experiments with information
displays, including testing and co-design with users, is a direction for
future research during ROBUST. The research described in this paper,
however, represents an initial exploration of this topic both through
case studies of current practices and use of a design-research approach.

The aim of this research was to improve understandings related to
mobility hub informational displays, and explore whether it was possible
to develop software tools for implementing mobility hubs at other
locations. The maps and other details shown in Figure
\ref{fig:display_outputs_figure}, and materials available online might
be suitable for prototyping, public consultation or similar\footnote{More
  formally, these results disprove the (null) hypothesis tested in this
  paper, by showing that it is possible to develop software to produce
  mobility hub information displays using OSM data.}. The flexibility of
R might also allow users to adjust layouts, styling and other details as
desired. Likewise, the OSM dataset can be freely accessed and edited by
anyone. Hence, starting with the developed software and a printer
practitioners, tactical urbanists or others might be able define a
mobility hub in their own community, and keep information displays up to
date as circumstances change.

This might raise larger questions, such as \emph{who gets to decide what
a mobility hub is?} Perhaps mobility hubs already exist wherever there
are combinations of: railway or electric vehicle charging stations; bus
or tram stops; shared bike docks, parking zones or other micro-mobility
offerings; or other transport-related facilities and services. Hopefully
the software tools developed in this study may help facilitate bottom-up
implementation, pop-ups or trials, or other efforts to develop and
maintain mobility hubs.

More generally, the research outcomes of this study may provide insights
about current mobility hub information displays developed by
practitioners, as a step towards identifying `best practices'. Mobility
hubs are still a new and evolving approach. As per the
\citet{Roukouni:2023aa} definition, a mobility hub needs to provide
``integration'', ``seamless, flexible connection'' and ``\ldots an
enjoyable experience for travelers'' across multiple modes, while also
connecting users to surrounding services and facilities. The research
report here, however, suggests that there may still be some way to go in
understanding how to deliver upon this definition, even within the
narrower confines of information displays and other guidance for users
at mobility hubs.

\renewcommand\refname{References}
\bibliography{References.bib, packages.bib}


\end{document}
